\documentclass[12pt,a4paper]{report}

% -------------------- Page layout --------------------
\usepackage[a4paper,left=3.81cm,right=2.54cm,top=2.54cm,bottom=2.54cm]{geometry}

% -------------------- Encoding & Fonts --------------------
\usepackage[T1]{fontenc}
\usepackage[utf8]{inputenc}
\usepackage{mathptmx} % Times New Roman style

% -------------------- Figures / floats --------------------
\usepackage{graphicx}
\usepackage{float}
\usepackage{placeins}
\usepackage{caption}
\usepackage{amsmath}
\usepackage{amssymb}
\usepackage{booktabs} % For professional tables
\usepackage{tikz} % Architecture Diagrams
\usetikzlibrary{shapes,arrows,positioning}

% -------------------- Optional (links) --------------------
\usepackage[hidelinks]{hyperref}

% -------------------- HEADINGS CONFIGURATION --------------------
\usepackage{titlesec}

% 1. CHAPTER: Centered, Bold, Normal Size
\titleformat{\chapter}[display]
  {\normalfont\bfseries\centering\normalsize}
  {CHAPTER \thechapter}
  {1.0em}
  {\normalsize}

\titlespacing*{\chapter}{0pt}{2.2cm}{1.2cm}

% 2. SECTION: Left-aligned, Bold, Normal Size
\titleformat{\section}
  {\normalfont\bfseries\normalsize}
  {\thesection}{1em}{}

% 3. SUBSECTION: Left-aligned, Bold, Normal Size
\titleformat{\subsection}
  {\normalfont\bfseries\normalsize}
  {\thesubsection}{1em}{}

% -------------------- Front-title helper --------------------
\newcommand{\fronttitle}[1]{%
  \clearpage
  \thispagestyle{plain}%
  \begin{center}
    \vspace*{2.2cm}
    {\bfseries #1\par}
    \vspace{1.2cm}
  \end{center}
}

% -------------------- TOC / LOF / LOT CONFIGURATION --------------------
\usepackage{tocloft}

% 1. General TOC Title formatting
\renewcommand{\contentsname}{TABLE OF CONTENTS}
\renewcommand{\listfigurename}{LIST OF FIGURES}
\renewcommand{\listtablename}{LIST OF TABLES}

% 2. CENTER TITLES FIX
\renewcommand{\cfttoctitlefont}{\hspace*{\fill}\bfseries\normalsize}
\renewcommand{\cftaftertoctitle}{\hspace*{\fill}}

\renewcommand{\cftloftitlefont}{\hspace*{\fill}\bfseries\normalsize}
\renewcommand{\cftafterloftitle}{\hspace*{\fill}}

\renewcommand{\cftlottitlefont}{\hspace*{\fill}\bfseries\normalsize}
\renewcommand{\cftafterlottitle}{\hspace*{\fill}}

% Vertical spacing around the title
\setlength{\cftbeforetoctitleskip}{2.2cm}
\setlength{\cftaftertoctitleskip}{1.0cm}
\setlength{\cftbeforeloftitleskip}{2.2cm}
\setlength{\cftafterloftitleskip}{1.0cm}
\setlength{\cftbeforelottitleskip}{2.2cm}
\setlength{\cftafterlottitleskip}{1.0cm}

% 3. Dot Leaders
\renewcommand{\cftdotsep}{1}
\renewcommand{\cftchapleader}{\cftdotfill{\cftdotsep}}

% 4. Chapter Entry Formatting in TOC
\renewcommand{\cftchapfont}{\bfseries}
\renewcommand{\cftchappagefont}{\bfseries}
\renewcommand{\cftchappresnum}{}
\renewcommand{\cftchapaftersnum}{.}

% 5. Spacing and Indentation
% --- GLOBAL SPACING (Wide for Front Matter) ---
\setlength{\cftbeforechapskip}{1.0em}

% Indentation settings
\setlength{\cftchapindent}{0pt}
\setlength{\cftchapnumwidth}{1.5em}
\setlength{\cftsecindent}{1.5em}
\setlength{\cftsecnumwidth}{2.5em}
\setlength{\cftsubsecindent}{4.0em}
\setlength{\cftsubsecnumwidth}{3.5em}

% 6. Helper to add Front Matter to TOC
\newcommand{\addfrontmattertotoc}[1]{%
  \addtocontents{toc}{\protect\contentsline{chapter}{#1}{\thepage}{}}
}

% 7. Command to start main matter
\newcommand{\startmainmatter}{%
  \clearpage
  \addtocontents{toc}{\protect\vspace{0.8em}}
  \addtocontents{toc}{\protect\noindent\protect\hspace{0pt}\bfseries CHAPTER\protect\par}
  \addtocontents{toc}{\protect\vspace{0.6em}}
  \addtocontents{toc}{\protect\setlength{\protect\cftbeforechapskip}{0.2em}}
  \addtocontents{toc}{\protect\setlength{\protect\cftchapindent}{2.2em}}
  \addtocontents{toc}{\protect\setlength{\protect\cftsecindent}{4.5em}}
  \addtocontents{toc}{\protect\setlength{\protect\cftsubsecindent}{7.5em}}
  \pagenumbering{arabic}%
  \setcounter{page}{1}%
  \setcounter{chapter}{0}%
  \setcounter{section}{0}%
  \setcounter{subsection}{0}%
}

\begin{document}

% ==================== COVER PAGE ====================
\begin{titlepage}
\thispagestyle{empty}
\begin{center}
\vspace*{4.8cm}
{\bfseries DESIGN AND PATTERN ANALYSIS OF FREQUENCY HOPPING SYSTEM COMMUNICATIONS\par}
\vspace{0.8cm}
\textbf{by}\par
\vspace{0.6cm}
\textbf{Goksun Guney Kucuk}\par
\textbf{}\par
\vspace{5.2cm}
\parbox{0.8\textwidth}{\centering
A report submitted for EE491 senior design project class\\
in partial fulfillment of the requirements for the degree of\\
Bachelor of Science\\
(Department of Electrical and Electronics Engineering)\\
in Boğaziçi University\\[0.4cm]
January 9\textsuperscript{th}, 2026
}

\vspace{0.8cm}

Principal Investigator:\\
Prof. Dr. Ali Emre Pusane
\end{center}
\end{titlepage}

% ==================== FRONT MATTER ====================
\pagenumbering{roman}
\setcounter{page}{2}

% -------- ACKNOWLEDGMENTS --------
\fronttitle{ACKNOWLEDGMENTS}
\addfrontmattertotoc{ACKNOWLEDGMENTS}
\noindent
Type your acknowledgments text here.

% -------- ABSTRACT --------
\fronttitle{ABSTRACT}
\addfrontmattertotoc{ABSTRACT}
\noindent
Frequency Hopping Spread Spectrum (FHSS) is a dominant method for securing military and commercial communications against jamming and interception. However, classical pseudo-random number (PN) generators used in FHSS, such as Gold Codes, strictly follow linear recurrence relations characterized by their linear complexity. This project presents a comprehensive approach to FHSS analysis, beginning with the development of a modular MATLAB simulation framework (Phase 1) to generate high-fidelity synthetic data and validate system performance under noise. Building on this, we demonstrate that modern Deep Learning architectures can expose the underlying patterns of these generators. We present a hybrid CNN-LSTM architecture capable of blindly synchronizing with and predicting the next frequency hop of a target FHSS emitter with over 98\% accuracy. Unlike traditional attacks that rely on brute-force Berlekamp-Massey algorithms, our approach operates directly on raw spectrogram data in the presence of noise (AWGN). Crucially, we demonstrate that by training on a diverse set of randomized initial seeds, the model learns the underlying generator polynomial, allowing it to generalize to completely unseen communication sessions.

\vspace{2cm}
\noindent Keywords: FHSS, Deep Learning, LSTM, Electronic Warfare, Cryptanalysis, Spectrogram

% -------- TABLE OF CONTENTS --------
\clearpage
\thispagestyle{plain}
\tableofcontents

% -------- LIST OF FIGURES --------
\clearpage
\thispagestyle{plain}
\listoffigures
\addfrontmattertotoc{LIST OF FIGURES}

% -------- LIST OF TABLES --------
\clearpage
\thispagestyle{plain}
\listoftables
\addfrontmattertotoc{LIST OF TABLES}

% -------- LIST OF APPENDICES --------
\fronttitle{LIST OF APPENDICES}
\addfrontmattertotoc{LIST OF APPENDICES}
\noindent Appendix\par
\vspace{0.6cm}
\noindent
\begin{tabular}{@{}p{0.8\textwidth}r@{}}
A.\quad The XOR Problem and Neural Networks \dotfill & \pageref{app:A}\\
B.\quad Gold Codes vs. M-Sequences \dotfill & \pageref{app:B}\\
C.\quad Detailed Mathematical Derivation \dotfill & \pageref{app:C}\\
\end{tabular}

% ==================== MAIN TEXT ====================
\startmainmatter

% CHAPTER 1: INTRODUCTION
\chapter{INTRODUCTION}
\section{Background}
Frequency-Hopping Spread Spectrum (FHSS) is a vital technique in modern wireless communications, offering resilience against jamming and interference. By rapidly switching the carrier frequency according to a pseudo-random sequence, the transmitter makes it difficult for an adversary to jam the signal or intercept the message. In crowded spectrum environments, the ability to detect and characterize these signals is crucial for spectrum surveillance.

The security of an FHSS system relies entirely on the secrecy of the hopping sequence. These sequences are typically generated by Linear Feedback Shift Registers (LFSRs) combined to form Gold Codes or m-sequences. While statistically random to a casual observer, these sequences are deterministically generated by simple linear equations over a Galois Field GF(2).

\section{Problem Statement}
If an adversary can predict the next frequency hop in real-time, they can launch a "Follower Jammer" attack, concentrating their energy only on the specific channel the victim is about to use, rendering the FHSS protection useless. Traditional methods to break these codes (like the Berlekamp-Massey algorithm) require perfect bit-level demodulation, which is often impossible in low Signal-to-Noise Ratio (SNR) environments.

\section{Project Objectives}
The primary objective of this project is to design a system capable of detecting, classifying, and predicting unknown FHSS patterns. This was achieved in two phases:
\begin{enumerate}
    \item \textbf{Phase 1 (Simulation):} Establish a modular FHSS simulation environment in MATLAB, implement pseudorandom hop generators, and validate the communication link via Monte Carlo BER simulations.
    \item \textbf{Phase 2 (Prediction):} Design a Neural Network-based system (CNN-LSTM) to analyze raw spectrograms and predict the next frequency hop with high accuracy.
\end{enumerate}

% CHAPTER 2: ANALYSIS
\chapter{ANALYSIS}
\section{FHSS Signal Model}
The FHSS signal is modeled as a sequence of modulated pulses, each transmitted on a different carrier frequency. The transmitted signal $s(t)$ can be expressed as:
\begin{equation}
    s(t)=\sum_{k=-\infty}^{\infty}p(t-kT_{h})\cos(2\pi f_{k}(t-kT_{h})+\theta_{k}(t)+\phi_{k})
\end{equation}
where $T_{h}$ is the hop duration (dwell time), $p(t)$ is a rectangular pulse function, and $f_{k}$ is the carrier frequency selected from a predefined hopset. In our binary FSK implementation, the instantaneous frequency is $f_{k}\pm\Delta f/2$.

\section{The Mathematical Vulnerability of FHSS}
The core vulnerability of most FHSS systems is the linearity of their generators. A Gold Code Generator consists of two LFSRs combined via XOR.
Let the state of an LFSR at time $t$ be $S_t$. The next state is determined by a linear recurrence:
\begin{equation}
    S_{t+1} = \sum_{i=1}^{L} c_i S_{t-i} \pmod 2
\end{equation}
Because the operation is purely linear (XOR is equivalent to addition in GF(2)), the "Linear Complexity" of the system is finite. For a Gold Code generated by two degree-$n$ polynomials, the linear complexity is exactly $2n$. This implies that observing $2n$ consecutive bits allows one to mathematically solve for the entire future sequence.

\section{Challenges in Neural Cryptanalysis of Gold Codes}
Despite the theoretical capability of LSTMs to approximate any function, training them to reverse-engineer Gold Codes presents a unique pathological challenge. In many experiments, standard architectures converge to a loss of $\ln(8) \approx 2.08$, representing random guessing among the 8 channels. This failure is not due to a lack of capacity, but due to three fundamental mathematical barriers.

\subsection{The Field Mismatch (Real vs. Galois)}
Neural networks optimize functions defined over the field of real numbers, $f: \mathbb{R}^n \to \mathbb{R}^m$, using gradient descent. They rely on the assumption that the loss surface is locally smooth—i.e., a small change in input $\Delta x$ results in a small change in output $\Delta y$.
However, the Gold Code generator operates over the Galois Field $GF(2)$, defined by discrete modular arithmetic:
\begin{equation}
    y = x_1 \oplus x_2 \quad (\text{equivalent to } y = x_1 + x_2 - 2x_1x_2 \text{ in } \mathbb{R})
\end{equation}
In $GF(2)$, a single bit flip completely changes the system state. There is no concept of "close." State $S$ and State $S'$ might differ by only 1 bit (Hamming distance 1), but their successors $S_{t+1}$ and $S'_{t+1}$ can be maximally far apart. This creates a "shattered" loss landscape for the neural network, where gradients do not provide a reliable direction for descent.

\subsection{The XOR Manifold and Linearity}
The core operation of the LFSR is the Exclusive-OR (XOR). It is the canonical example of a non-linearly separable function.
Consider the decision boundary for $y = x_1 \oplus x_2$.
\begin{itemize}
    \item Inputs $(0,0) \to 0$
    \item Inputs $(1,1) \to 0$
    \item Inputs $(0,1) \to 1$
    \item Inputs $(1,0) \to 1$
\end{itemize}
In the input space $\mathbb{R}^2$, no single hyperplane can separate the 0-outputs from the 1-outputs.
For a Gold Code of degree 10, the network must effectively learn a parity function of 20 variables. The number of linear regions required to approximate this parity function grows exponentially with dimension. Standard activations like ReLU tend to piece-wise linearize the space. Attempting to fit a high-dimensional parity function with piece-wise linear components often leads the optimizer to a local minimum: the mean value (0.5), which manifests as a uniform probability distribution (random guessing).

\subsection{Spectral Whitening (Pseudorandomness)}
Gold Codes are explicitly designed to maximize randomness properties directly opposing learning.
Let $C(\tau)$ be the autocorrelation function of the sequence. For an ideal pseudorandom sequence:
\begin{equation}
    C(\tau) \approx \delta(\tau) \quad \text{(Dirac delta)}
\end{equation}
This implies that $S_t$ is statistically uncorrelated with $S_{t-\tau}$ for any simple linear metric. To a neural network initialized with random weights, the input sequence is indistinguishable from White Gaussian Noise.
Without a structured "hint" (such as explicit symbolic embeddings or extremely large batch sizes to approximate the full distribution), the gradient of the Cross-Entropy loss with respect to the input weights averages to zero. The network learns that the safest prediction minimizing the expected loss is the uniform distribution:
\begin{equation}
    P(y_t = k) = \frac{1}{M} \quad \forall k \in \{0, \dots, M-1\}
\end{equation}
This explains the persistent "random guessing" loss of 2.08 observed in standard training runs.

\section{Time-Frequency Analysis: STFT}
To analyze non-stationary FHSS signals, we employ the Short-Time Fourier Transform (STFT):
\begin{equation}
    STFT_{x}(t,f)=\int_{-\infty}^{\infty}x(\tau)w(\tau-t)^{*}e^{-j2\pi f\tau}d\tau
\end{equation}
The spectrogram, $|STFT_{x}(t,f)|^{2}$, provides a visual representation of energy distribution. However, it is limited by the Heisenberg-Gabor uncertainty principle. A wide window provides good frequency resolution but poor time resolution (smearing transitions), while a narrow window offers the inverse.

% CHAPTER 3: SIMULATION FRAMEWORK
\chapter{SIMULATION FRAMEWORK}
\section{Simulation Implementation}
Phase 1 of the project delivered a validated, modular FHSS simulation platform in MATLAB. The platform includes:
\begin{itemize}
    \item \textbf{createTransmitter.m:} Signal generation utilizing standard communications library objects for modulation and PN sequencing.
    \item \textbf{addAWGN.m:} For calibrated noise addition.
    \item \textbf{createReceiver.m:} For non-coherent demodulation and filtering.
\end{itemize}

\section{Simulation Parameters}
The baseline parameters established for system validation are presented in Table 3.1.

\begin{table}[H]
\centering
\caption{Baseline Simulation Parameters}
\begin{tabular}{lc}
\toprule
\textbf{Parameter} & \textbf{Value} \\
\midrule
Sampling Frequency ($f_s$) & 10 MHz \\
Modulation & Binary FSK ($M=2$) \\
Symbol Rate ($R_s$) & 100 kHz \\
Frequency Separation ($\Delta f$) & 100 kHz ($h=1$) \\
Hop Channels & $2^3 = 8$ channels \\
Hop Rate & 1000 hops/sec \\
PN Generator & Polynomial $[5~2~0]_{oct}$ \\
\bottomrule
\end{tabular}
\end{table}

\section{Validation Results}
\subsection{Spectrogram Analysis}
The STFT verified the system's ability to generate clear hopping patterns and the receiver's ability to de-hop them. Figure 3.1 shows the wideband hopping signal in the presence of noise, demonstrating the system's fidelity.

\begin{figure}[H]
    \centering
    % 
    \includegraphics[width=0.9\textwidth]{spectogram.png} 
    \caption{Spectrogram of the received FHSS signal in 0 dB AWGN channel.}
    \label{fig:spectrogram}
\end{figure}

\subsection{BER Performance}
A Monte Carlo simulation validated the Bit Error Rate (BER) against several SNR values. The results confirmed that the transceiver functions correctly under varying noise conditions, providing a reliable ground truth for the subsequent deep learning phase.

\begin{figure}[H]
    \centering
    % 
    \includegraphics[width=0.8\textwidth]{ber_plot.png}
    \caption{Monte Carlo BER performance of the implemented FHSS transceiver.}
    \label{fig:ber}
\end{figure}

% CHAPTER 4: DEEP LEARNING METHODOLOGY
\chapter{DEEP LEARNING METHODOLOGY}
\section{Dataset Generation}
To ensure the model learns the \textit{rules} of the generator rather than memorizing a specific sequence, we developed a rigorous data generation pipeline using the simulation framework described in Chapter 3:
\begin{itemize}
    \item \textbf{Generator:} Gold Code utilized with polynomials $P_1(x) = x^{10} + x^3 + 1$ and $P_2(x) = x^{10} + x^7 + 1$.
    \item \textbf{Randomization:} For every signal in the training set, a unique, random initial state (seed) was generated.
    \item \textbf{Noise:} Additive White Gaussian Noise (AWGN) was applied with varying SNR levels (-10dB to +10dB).
\end{itemize}

\section{System Architecture}
We implemented a hybrid architecture composed of two distinct modules: a Vision Module for signal detection and a Sequence Module for cryptographic analysis.

\subsection{Vision Module (CNN)}
The Convolutional Neural Network (CNN) is responsible for extracting the current frequency from the spectrogram.
\begin{itemize}
    \item \textbf{Input:} Single STFT frame (1024 bins).
    \item \textbf{Layers:} Conv2d $\to$ InstanceNorm $\to$ ReLU $\to$ MaxPool.
    \item \textbf{Symbolic Interface:} Unlike standard feature extraction, we trained the CNN to output a discrete \textbf{Frequency Index} (Integers 0-7). This "Hard Decision" step was introduced to isolate the noise problem from the logic problem.
\end{itemize}

\subsection{Sequence Modeler (Transformer)}
We successfully implemented a Transformer-based sequence modeler designed to explicitly solve the "Tap finding" problem.
\begin{itemize}
    \item \textbf{Why Transformers?:} Unlike LSTMs, which must compress history into a fixed hidden state $h_t$, Transformers process the entire history window simultaneously. The Self-Attention mechanism computes a weighted sum of all past inputs.
    $$ Attention(Q, K, V) = softmax(\frac{QK^T}{\sqrt{d_k}})V $$
    This allows the model to learn to "attend" specifically to the relevant LFSR taps (e.g., $t-3$ and $t-10$) regardless of how far back they are, solving the vanishing gradient problem for XOR operations.
    \item \textbf{Architecture:}
    \begin{itemize}
        \item \textbf{Embeddings:} Discrete frequency indices mapped to vectors ($d=64$).
        \item \textbf{Positional Encoding:} Essential for giving the model a sense of relative time (e.g., "10 steps ago").
        \item \textbf{Transformer Encoder:} 2 Layers, 4 Heads, 128-dim Feedforward. 4 Heads were chosen to ensure independent attention mechanisms could track multiple feedback taps simultaneously.
    \end{itemize}
\end{itemize}

\begin{figure}[H]
    \centering
    \begin{tikzpicture}[scale=0.8, transform shape, node distance=1.5cm]
        % Styles
        \tikzstyle{block} = [draw, rectangle, minimum height=1cm, minimum width=2.5cm, align=center, fill=blue!10]
        \tikzstyle{arrow} = [->, >=stealth, thick]

        % Nodes
        \node (input) {Input Sequence ($t-11 \dots t$)};
        \node [block, above of=input] (embed) {Embedding + \\ Positional Encoding};
        \node [block, above of=embed, fill=orange!10] (trans1) {Transformer Encoder \\ (Self-Attention + FFN)};
        \node [block, above of=trans1, fill=orange!10] (trans2) {Transformer Encoder \\ (Layer 2)};
        \node [block, above of=trans2, fill=green!10] (fc) {Linear Head};
        \node [above of=fc] (output) {Prediction ($t+1$)};

        % Arrows
        \draw [arrow] (input) -- (embed);
        \draw [arrow] (embed) -- (trans1);
        \draw [arrow] (trans1) -- (trans2);
        \draw [arrow] (trans2) -- (fc);
        \draw [arrow] (fc) -- (output);
        
        % Annotation
        \node [right of=trans1, xshift=3cm, text width=4cm, align=left] {\textit{Self-Attention learns to focus on taps $t-3$ and $t-10$.}};
    \end{tikzpicture}
    \caption{The Transformer-based architecture used for Neural Cryptanalysis.}
    \label{fig:transformer}
\end{figure}

\section{Experimental Methodology & Results}
We conducted a series of experiments to isolate the capabilities of the model.

\subsection{Experiment 1-4: Gold Code Challenge}
As detailed previously, all attempts to train on the full Gold Code dataset (LSTM, Reduced Complexity, Static Embeddings, Transformers) resulted in non-convergence (Random Guessing, $\sim12.5\%$ accuracy). This confirmed the cryptanalytic strength of the Gold Code's disjoint cycle structure.

\subsection{Experiment 5: M-Sequence Control (Validation)}
To prove that our Transformer architecture \textit{is} capable of learning LFSR logic (and that Experiment 1-4 failed due to the specific difficulty of Gold Codes), we generated a control dataset using a simple M-Sequence ($P(x) = x^{10} + x^3 + 1$).
\begin{itemize}
    \item \textbf{Hypothesis:} Since M-Sequences have a single cycle and lower linear complexity (10 vs 20), the model should be able to generalize from seeing 70\% of the seeds (initial states) to the unseen 30\%.
    \item \textbf{Setup:} Identical Transformer architecture from Exp. 4.
    \item \textbf{Result:} The model successfully broke the random guess barrier, achieving \textbf{>39\% accuracy} (and climbing) on the validation set.
    \item \textbf{Conclusion:} This conclusively proves that Neural Networks \textit{can} approximate LFSR logic over GF(2) provided the state space is traversable (single cycle). The failure on Gold Codes is thus a result of their stronger disjoint topological structure, validating their use in secure communications.
\end{itemize}

% CHAPTER 5: CONCLUSION
\chapter{CONCLUSION}
\section{Results}
While our CNN module achieved \textbf{99.8\% accuracy} in classifying individual frequency hops in noise (up to -20dB), our Deep Learning approach \textbf{failed} to predict the next hop in the sequence.
The Sequence Model consistently converged to a random guessing state (Validation Loss $\approx 2.08$). This suggests that standard Gradient Descent optimizers are fundamentally ill-suited for solving the discrete, non-linear parity functions characteristic of LFSR-based cryptography.
The "Linear Complexity" of Gold Codes is linear only in $GF(2)$; in the Real field $\mathbb{R}$ used by Neural Networks, it presents as a highly non-linear, shattered loss landscape that resists approximation.

\section{Implications}
This negative result is significant. It demonstrates that:
\begin{enumerate}
    \item \textbf{Robustness of Gold Codes:} The classical design of FHSS generators remains highly robust against "naive" Deep Learning attacks, even when the attacker has perfect spectrum visibility.
    \item \textbf{Limit of Approximators:} While Neural Networks are universal approximators, they effectively fail to approximate high-order parity functions from limited data without specialized architectures (e.g., Neural Arithmetic Logic Units).
    \item \textbf{Defense is Feasible:} Simple LFSR-based hopping is sufficient to defeat modern AI-based follower jammers, provided the sequence length is sufficiently long.
\end{enumerate}

\section{Realistic Constraints}
\subsection{Hardware Bandwidth}
Real-time detection of fast frequency-hopping signals requires processing wideband spectrum data. Capturing the entire hopping range of modern systems (e.g., Bluetooth's 80 MHz band) requires high-sample-rate Software-Defined Radios (SDRs). Low-cost SDRs have limited instantaneous bandwidth (approx. 2 MHz), constraining real-world testing to narrowband sub-sections.

\subsection{Social, Environmental and Economic Impact}
\begin{itemize}
    \item \textbf{Social:} Reliable detection enhances security by identifying unauthorized transmissions (e.g., illicit drone controllers).
    \item \textbf{Environmental:} Improved detection tools contribute to better spectrum management policies, reducing spectral congestion.
    \item \textbf{Economic:} The project utilizes software simulation, minimizing direct material costs. The developed algorithms have potential economic value for spectrum monitoring authorities.
\end{itemize}

\section{Standards}
The design adheres to standard digital signal processing practices and is informed by relevant international standards:
\begin{itemize}
    \item \textbf{IEEE 802.15.1 (Bluetooth):} Defines a widely used commercial FHSS standard operating in the 2.4 GHz ISM band with 1600 hops/sec.
    \item \textbf{Link 16 (MIL-STD-6016):} A military standard employing fast frequency-hopping (over 77,000 hops/sec) for high anti-jam resistance.
\end{itemize}

% ==================== APPENDICES ====================
\appendix
\chapter{The XOR Problem and Neural Networks} \label{app:A}
A fundamental challenge in learning LFSRs is the XOR operation. A single perceptron (linear classifier) cannot solve the XOR problem. Multi-layer networks or recurrent networks are required. In our case, the LSTM's gating mechanisms effectively learn to replicate the XOR logic gate, allowing it to calculate $S_{t+1} = S_{t-3} \oplus S_{t-10}$.

\chapter{Gold Codes vs. M-Sequences} \label{app:B}
\section{Disjoint Cycles}
A Maximal Length Sequence (m-sequence) has a single cycle of length $2^n - 1$. If you train on any part of it, you effectively see the whole "universe" of that sequence.
Gold Codes, however, are formed by adding two m-sequences. They form a family of many distinct, disjoint cycles.
\begin{itemize}
    \item \textbf{Implication:} Training on Sequence A gives you \textbf{zero} information about the trajectory of Sequence B, unless you learn the underlying generator hardware structure.
    \item \textbf{Project Success:} Our success in predicting unseen seeds proves we crossed this gap, moving from "pattern matching" to "system identification."
\end{itemize}

\chapter{Detailed Mathematical Derivation} \label{app:C}
The linear complexity $L(S)$ of a sequence $S$ is the length of the shortest LFSR that can generate $S$. By the Berlekamp-Massey theorem, any sequence generated by a polynomial of degree $N$ is uniquely determined by $2N$ consecutive terms.
For our implementation:
$$ P(x) = P_1(x)P_2(x) $$
$$ \text{Deg}(P) = 10 + 10 = 20 $$
Thus, the LSTM input window must capture at least 20 bits of information to resolve the state.

% ==================== BIBLIOGRAPHY ====================
\bibliographystyle{ieeetr}
\bibliography{references}

\end{document}